\begin{figure}[htbp]
	\centering
	\begin{tikzpicture}[font=\small]
		% ---------------------------------------------------------------
		% (a) trave come corpo: prisma in assonometria e linea d'asse
		% ---------------------------------------------------------------
		\begin{scope}
			\def\L{4.2}\def\H{0.55}\def\dx{0.40}\def\dy{0.28}
			% spigoli nascosti
			\draw[trvcostruzione] (0,0) -- (\dx,\dy) -- (\L+\dx,\dy);
			\draw[trvcostruzione] (\dx,\dy) -- (\dx,\H+\dy);
			% faccia frontale, superiore, laterale destra
			\draw[trvconcio] (0,0) rectangle (\L,\H);
			\draw[trvconcio] (0,\H) -- (\dx,\H+\dy) -- (\L+\dx,\H+\dy) -- (\L,\H);
			\draw[trvconcio] (\L,0) -- (\L+\dx,\dy) -- (\L+\dx,\H+\dy);
			% linea d'asse (baricentri delle sezioni)
			\draw[line width=.6pt, dash pattern=on 6pt off 1.5pt on 1pt off 1.5pt]
				(-0.35+\dx/2,\H/2+\dy/2) -- (\L+\dx/2+0.35,\H/2+\dy/2);
			\pic at (\dx/2,\H/2+\dy/2) {trvpunto};
			\pic at (\L+\dx/2,\H/2+\dy/2) {trvpunto};
			\draw[trvrichiamo] (\dx/2,\H/2+\dy/2) -- (\dx/2-0.25,\H+\dy+0.3) node[above, black] {$A$};
			\draw[trvrichiamo] (\L+\dx/2,\H/2+\dy/2) -- (\L+\dx/2+0.25,\H+\dy+0.3) node[above, black] {$B$};
			% quote
			\draw[trvquota] (0,-0.35) -- node[below=1pt] {$\ell$} (\L,-0.35);
			\draw[trvquota] (-0.55,0) -- node[left=1pt] {$h$} (-0.55,\H);
			\draw[trvquota] (\L+0.25,-0.05) -- node[below right=-1pt] {$b$} (\L+0.25+\dx,\dy-0.05);
			\node at (\L/2,-1.15) {(a)};
		\end{scope}
		% ---------------------------------------------------------------
		% (b) schematizzazione: linea d'asse e vincoli
		% ---------------------------------------------------------------
		\begin{scope}[xshift=5.9cm, yshift=0.42cm]
			\def\L{4.2}
			\draw[trvtrave] (0,0) -- (\L,0);
			\pic at (0,0) {polcerniera};
			\pic at (\L,0) {polcarrello};
			\node[above=2pt] at (0,0) {$A$};
			\node[above=2pt] at (\L,0) {$B$};
			\draw[trvquota] (0,-0.77) -- node[below=1pt] {$\ell$} (\L,-0.77);
			\node at (\L/2,-1.57) {(b)};
		\end{scope}
	\end{tikzpicture}
	\caption{Dalla trave alla schematizzazione strutturale.
		(a)~Trave di lunghezza $\ell$ e dimensioni
		trasversali $b$, $h \ll \ell$; la linea d'asse,
		di estremi $A$ e $B$, è il luogo dei baricentri
		delle sezioni.
		(b)~Schematizzazione: si traccia la sola linea
		d'asse, con i vincoli (qui cerniera in $A$ e
		carrello in $B$) nei punti corrispondenti.}
	\label{fig:trave_convenzione}
\end{figure}
